% T9 -- Paired bootstrap over the test users (slide 75).
\begin{table}[!htbp]
\centering\small
\caption{Is any step up the ladder real? Every interval says no}
\label{tab:bootstrap}
\begin{tabular}{@{}lrlc@{}}
\toprule
\textbf{Step in training people} & \textbf{Measured difference} & \textbf{Range it could really be} & \textbf{Real gain?} \\
\midrule
5{,}000 $\rightarrow$ 10{,}000  & $+0.20\pt$ & $-0.21 \ \ldots \ +0.64$ & no \\
10{,}000 $\rightarrow$ 15{,}000 & $+0.09\pt$ & $-0.32 \ \ldots \ +0.52$ & no \\
15{,}000 $\rightarrow$ 21{,}000 & $-0.25\pt$ & $-0.72 \ \ldots \ +0.20$ & no \\
\midrule
5{,}000 $\rightarrow$ 21{,}000 \textit{(total)} & $+0.04\pt$ & $-0.43 \ \ldots \ +0.51$ & no \\
\bottomrule
\end{tabular}

\figtext{%
\figpara{The question.}{Between two training sizes the score moves a little. Is
that movement the model genuinely improving, or is it the luck of which thousand
people happened to be in the test set?}

\figpara{How it is tested.}{A thousand people are drawn at random from the
thousand held-out ones, with replacement, so some appear twice and others not at
all. Both models are scored on exactly that draw and the difference is noted. That
is repeated two thousand times, and the middle 95\,\% of the two thousand
differences is the range printed in the third column. Because both models always
see the same draw, the luck of the sample cancels and only the difference between
the models is left.}

\figpara{How to read a row.}{If the range contains zero, then ``no difference at
all'' is a perfectly ordinary outcome of this experiment, and the two sizes
cannot be told apart. Every row here contains zero, including the last one, which
spans the entire ladder from 5{,}000 to 21{,}000 people and whose best estimate
is 0.04.}

\figpara{What that licenses me to say.}{Not ``the gain from more people is
small''. \emph{No gain was measured at all} above 5{,}000 people. \seleven,
1{,}000 test people, 2{,}000 resamples.}}

\vspace{3pt}
\begin{minipage}{\textwidth}\footnotesize
The fitted scaling law says the same thing: a slope of $+0.18\pt$ per ten-fold
increase in users, and a line that accounts for only 12\,\% of the spread of
the eight points it is fitted to. At that slope one extra point would require
about 1.5 billion users, the absurdity that shows the curve is flat rather
than merely slow. The correct statement is not ``the gain is small'' but
``no gain was measured''.
\end{minipage}
\figsource{\nbfinal}{train\_cellid\_final.ipynb}{the paired-bootstrap cells, 2,000 resamples}
\end{table}
